\documentclass[UTF8,zihao=-4]{ctexart}
\usepackage[a4paper,margin=1.9cm]{geometry}
\usepackage{hyperref}
\usepackage{array}
\usepackage{longtable}
\usepackage{verbatim}
\usepackage{xcolor}
\hypersetup{hidelinks}
\setlength{\parindent}{0pt}
\setlength{\parskip}{0.45em}
\emergencystretch=4em
\sloppy
\title{模拟卷08-2025风格题目与详解}
\author{}
\date{}
\begin{document}
\maketitle
\setcounter{tocdepth}{1}
\tableofcontents
\newpage
\section{编译原理 2025 风格模拟卷08（题目与详解）}


\subsection{卷面说明}


本卷按 2025 年期末卷组织，主线题为词法分析、LL、指令调度、支配、符号执行、Andersen。最后设置旧卷加练题，覆盖 2022、2023 中仍有价值的传统题型。总分 100 分，加练题 10 分。


\section{A 卷题目}


\subsection{一、词法分析（15 分）}


正则表达式：


\begin{quote}
\footnotesize
\begin{verbatim}
a* b (a|b)*
\end{verbatim}
\end{quote}


\begin{enumerate}
\item 使用 Thompson 构造法构造等价 NFA。
\item 使用子集构造法构造 DFA，写明 DFA 状态含义。
\item 最小化 DFA。
\end{enumerate}


\subsection{二、LL 语法分析（15 分）}


文法：


\begin{quote}
\footnotesize
\begin{verbatim}
K -> L M
L -> L hash | key
M -> open | close
\end{verbatim}
\end{quote}


\begin{enumerate}
\item 消除左递归。
\item 写 FIRST 和 FOLLOW。
\item 写预测分析表。
\item 判断消除左递归后的文法是否为 LL(1)。
\end{enumerate}


\subsection{三、指令调度（15 分）}


给定指令：


\begin{quote}
\footnotesize
\begin{verbatim}
I1: R1 = R2 + R3
I2: R4 = R1 + R5
I3: R2 = R6 + R7
I4: R1 = R8 + R9
I5: R10 = R4 + R1
\end{verbatim}
\end{quote}


\begin{enumerate}
\item 写出指令调度的目的。
\item 写出 RAW、WAR、WAW 的定义。
\item 构建依赖图，WAR/WAW 用虚线表示。
\item 在允许寄存器重命名的条件下给出一个合法重排序列。
\end{enumerate}


\subsection{四、支配问题（15 分）}


CFG：循环带出口分支


\begin{quote}
\footnotesize
\begin{verbatim}
Entry->H
H->Body
Body->If
If->Body
If->Exit
\end{verbatim}
\end{quote}


\begin{enumerate}
\item 写出 Dom、PostDom、Dom Frontier 的定义。
\item 计算每个节点的 Dom 集和 idom。
\item 写出支配边界。
\item 证明支配关系图是森林。
\end{enumerate}


\subsection{五、符号执行与 SAT（20 分）}


公式：


\begin{quote}
\footnotesize
\begin{verbatim}
F = ((a8 and b8) or c8)
\end{verbatim}
\end{quote}


\begin{enumerate}
\item 写出 Path-sensitive、Flow-sensitive、Context-sensitive 的定义。
\item 写出 CNF 的定义。
\item 使用 Tseitin 算法把 F 转成 CNF。
\item 证明 2025 卷中的两组公式同时可满足。
\end{enumerate}


\subsection{六、Andersen 指针分析（20 分）}


程序：


\begin{quote}
\footnotesize
\begin{verbatim}
a = &O1
b = &O2
c = a
d = b
*c = d
e = *a
\end{verbatim}
\end{quote}


\begin{enumerate}
\item 对每条语句写 Andersen 约束。
\item 根据约束构建闭包。
\item 写出最终 points-to 集合。
\end{enumerate}


\subsection{七、旧卷加练题（10 分）}


类型：LR


说明 LR(1) 项目 closure、goto 和归约设置条件。


\section{B 详解}


\subsection{一、词法分析详解}


\begin{quote}
\footnotesize
\begin{verbatim}
正则表达式：a* b (a|b)*

Thompson NFA：起点 q0，终点 q7。按下列转移表画图，q7 画双圈。
q2 --a--> q3
q0 --epsilon--> q2
q0 --epsilon--> q1
q3 --epsilon--> q2
q3 --epsilon--> q1
q4 --b--> q5
q1 --epsilon--> q4
q10 --a--> q11
q12 --b--> q13
q8 --epsilon--> q10
q8 --epsilon--> q12
q11 --epsilon--> q9
q13 --epsilon--> q9
q6 --epsilon--> q8
q6 --epsilon--> q7
q9 --epsilon--> q8
q9 --epsilon--> q7
q5 --epsilon--> q6

子集构造 DFA：
D0:
  NFA集合 = {q0,q1,q2,q4}
  on a -> D1
  on b -> D2
  终态 = 否
D1:
  NFA集合 = {q1,q2,q3,q4}
  on a -> D1
  on b -> D2
  终态 = 否
D2:
  NFA集合 = {q5,q6,q7,q8,q10,q12}
  on a -> D3
  on b -> D4
  终态 = 是
D3:
  NFA集合 = {q7,q8,q9,q10,q11,q12}
  on a -> D3
  on b -> D4
  终态 = 是
D4:
  NFA集合 = {q7,q8,q9,q10,q12,q13}
  on a -> D3
  on b -> D4
  终态 = 是

最小化：
初分组：
  终态组 = {D2,D3,D4}
  非终态组 = {D0,D1}
最终分组：
  G1 = {D2,D4,D3}
  G2 = {D0,D1}
每个最终分组对应一个最小 DFA 状态。
\end{verbatim}
\end{quote}



\subsection{二、LL 语法分析详解}


\begin{quote}
\footnotesize
\begin{verbatim}
原文法：
K -> L M
L -> L hash | key
M -> open | close

消除左递归：
K  -> L M
L  -> key L'
L' -> hash L' | epsilon
M  -> open | close

FIRST：
FIRST(K)={ key }
FIRST(L)={ key }
FIRST(L')={ hash, epsilon }
FIRST(M)={ open, close }

FOLLOW：
FOLLOW(K)={ $ }
FOLLOW(L)={ open, close }
FOLLOW(L')={ open, close }
FOLLOW(M)={ $ }

预测表：
M[K,key] = K -> L M
M[L,key] = L -> key L'
M[L',hash] = L' -> hash L'
M[L',open] = L' -> epsilon
M[L',close] = L' -> epsilon
M[M,open] = M -> open
M[M,close] = M -> close

结论：表中无多重入口，消除左递归后的文法是 LL(1)。
\end{verbatim}
\end{quote}



\subsection{三、指令调度详解}


定义：


\begin{quote}
\footnotesize
\begin{verbatim}
指令调度在不改变程序语义的前提下重排指令，减少流水线停顿，提高并行执行效率。
RAW：先写后读，必须保持顺序。
WAR：先读后写，写者不能提前。
WAW：先写后写，顺序影响最终值。
\end{verbatim}
\end{quote}


\begin{quote}
\footnotesize
\begin{verbatim}
指令：
I1: R1 = R2 + R3
I2: R4 = R1 + R5
I3: R2 = R6 + R7
I4: R1 = R8 + R9
I5: R10 = R4 + R1

依赖：
I1 -> I2 RAW on R1
I1 -> I3 WAR on R2
I1 -> I4 WAW on R1
I2 -> I4 WAR on R1
I2 -> I5 RAW on R4
I4 -> I5 RAW on R1

寄存器重命名：
I3: R2_new = R6 + R7
I4: R1_new = R8 + R9
I5 使用 R1_new

重命名后保留 RAW，WAR/WAW 被消除。一个合法调度序列：
I1, I3, I4, I2, I5
\end{verbatim}
\end{quote}



\subsection{四、支配问题详解}


定义：


\begin{quote}
\footnotesize
\begin{verbatim}
Dom(n)：入口到 n 的每条路径都经过的节点集合。
PostDom(n)：n 到出口的每条路径都经过的节点集合。
Dom Frontier(n)：n 支配某个前驱，但不严格支配该节点的位置。
\end{verbatim}
\end{quote}


\begin{quote}
\footnotesize
\begin{verbatim}
CFG 边：
Entry->H
H->Body
Body->If
If->Body
If->Exit

Dom 集：
Dom(Entry)={Entry}
Dom(H)={Entry,H}
Dom(Body)={Entry,H,Body}
Dom(If)={Entry,H,Body,If}
Dom(Exit)={Entry,H,Body,If,Exit}

直接支配者：
idom(H)=Entry
idom(Body)=H
idom(If)=Body
idom(Exit)=If

支配边界：
DF(Body)={Body}
DF(If)={Body}
其余节点 DF 为空集

森林证明：除入口外，每个可达节点有唯一直接支配者；直接支配关系严格支配且无环；所以每个节点至多一个父节点，整体构成树，多入口时构成森林。
\end{verbatim}
\end{quote}



\subsection{五、符号执行与 SAT 详解}


\begin{quote}
\footnotesize
\begin{verbatim}
公式：F = ((a8 and b8) or c8)
引入：x81 <-> (a8 and b8)，x82 <-> (x81 or c8)，根变量 x82 为真。

CNF：
(not x81 or a8)
(not x81 or b8)
(x81 or not a8 or not b8)
(x82 or not x81)
(x82 or not c8)
(not x82 or x81 or c8)
(x82)

敏感性定义：
Path-sensitive 区分不同路径。
Flow-sensitive 区分程序点和语句顺序。
Context-sensitive 区分函数调用上下文。

同时可满足证明套用：
F1=(a0 or ... or an or c) and (b0 or ... or bm or not c)，F2=(a0 or ... or an or b0 or ... or bm)。
F1 为真时，c=false 推出某个 ai=true，c=true 推出某个 bj=true，所以 F2 为真。
F2 为真时，若某个 ai=true 令 c=false；若某个 bj=true 令 c=true，所以 F1 为真。
\end{verbatim}
\end{quote}



\subsection{六、Andersen 指针分析详解}


\begin{quote}
\footnotesize
\begin{verbatim}
程序：
a = &O1
b = &O2
c = a
d = b
*c = d
e = *a

约束：
O1 in pts(a)
O2 in pts(b)
pts(a) subset pts(c)
pts(b) subset pts(d)
for v in pts(c), pts(d) subset pts(v)
for v in pts(a), pts(v) subset pts(e)

闭包：
pts(a)={O1}
pts(b)={O2}
pts(c)={O1}
pts(d)={O2}
由 *c=d 得 pts(O1)={O2}
由 e=*a 得 pts(e)={O2}
\end{verbatim}
\end{quote}



\subsection{七、旧卷加练题详解}


\begin{quote}
\footnotesize
\begin{verbatim}
closure：若 [A -> alpha . B beta, a] 在 I 中，对 B->gamma 和 b in FIRST(beta a)，加入 [B -> . gamma, b]。goto(I,X)=closure({[A -> alpha X . beta, a]})。归约项 [A -> alpha ., a] 只在 a 列填 reduce。
\end{verbatim}
\end{quote}


\section{C 复盘清单}


\begin{quote}
\footnotesize
\begin{verbatim}
词法：Thompson、epsilon-closure、move、终态集合、最小化分组。
LL：左递归、FIRST、FOLLOW、预测表、多重入口。
调度：RAW/WAR/WAW、重命名、拓扑序。
支配：Dom、idom、DF、森林证明。
SAT：敏感性、CNF、Tseitin 子句、同时可满足证明。
Andersen：四类约束、闭包传播、最终 points-to 集。
\end{verbatim}
\end{quote}

\end{document}
